% Generated draft: requires mathematical and source review.
\documentclass[11pt]{article}
\usepackage[margin=25mm]{geometry}
\usepackage{fontspec}
\setmainfont{Latin Modern Roman}
\usepackage{amsmath,amssymb,mathtools,unicode-math}
\setmathfont{Latin Modern Math}
\usepackage{enumitem,needspace,xurl,hyperref}
\setlength{\parindent}{0pt}
\setlength{\parskip}{5pt}
\setlength{\emergencystretch}{3em}
\newcommand{\N}{\mathbb{N}}
\newcommand{\Z}{\mathbb{Z}}
\newcommand{\Q}{\mathbb{Q}}
\newcommand{\R}{\mathbb{R}}
\newcommand{\C}{\mathbb{C}}
\newcommand{\F}{\mathbb{F}}
\newcommand{\A}{\mathbb{A}}
\newcommand{\PP}{\mathbb{P}}
\newcommand{\E}{\mathbb{E}}
\newcommand{\eps}{\varepsilon}
\newcommand{\dd}{\,\mathrm{d}}
\begin{document}
\section*{UnsolvedMath ID: 10400123}
\textbf{Source problem code:} catalogue entry
\subsection{Definitions and mathematical statement}
The source asks about ``such an expansion'', referring to omitted preceding material. Retain its three requests: establish existence of the indicated expansion; understand its relationship with Ohtsuki's and Habiro's expansions; and give independent topological meaning to its resulting invariants, especially its phases \(c_j\).

The highlighted phase-identification target is that these phases are the corresponding Chern--Simons values,
\[
 c_j=\operatorname{CS}(A_j)
\]
in the original normalization and period convention, with the relevant connections or moduli components \(A_j\) still to be specified. The excerpt supplies neither the expanded invariant, expansion parameter and form, coefficient definitions, manifold class, gauge group nor indexing of phases. No particular oscillatory asymptotic formula or power-series completion is reconstructed from the word ``expansion''.

The passage ends by announcing ``the following conjecture'' but does not include that conjecture. Its semi-classical discussion is therefore incomplete. The summary separately mentions a measure \(\tau_M(A)^{1/2}\), but gives no definition or precise assertion; that fragment cannot stand in for the missing expansion or conjecture. The cited Ohtsuki/Habiro relationships likewise need the exact input classes and comparison maps.

The supplied open/triage label contains no resolving theorem. These source requests are retained as an incomplete research program, not certified as a well-posed universal expansion conjecture or independently verified current status.
\subsection{Short English statement}
Establish the intended expansion, relate it to Ohtsuki and Habiro expansions, and identify its phases with Chern--Simons data. The expansion itself and the announced next conjecture are missing, so the exact invariant and phase statement cannot be recovered.
\subsection{Sources}
\begin{itemize}
\item Local registry: \url{data/unsolvedmath/problems.json}, permanent UnsolvedMath ID \(10400123\), source code catalogue entry; supplied statement, background, and research summary, updated 2026-08-06.
\item Ohtsuki, \emph{Problems on invariants of knots and 3-manifolds} (2002), Problem 7.8, attributed to J. E. Andersen, PDF page 106: \url{https://msp.org/gtm/2002/04/gtm-2002-04-024s.pdf}. Supplied PDF-font-structure extraction accessed 2026-07-29.
\item Ohtsuki and Habiro expansions and the summary's \(\tau_M(A)^{1/2}\) measure are mentioned without separate definitions or bibliographic entries; the announced following conjecture is absent.
\end{itemize}
\end{document}
